% Source: https://artofproblemsolving.com/community/c3h1548800_circles_with_centers_at_the_vertices_of_regular_hexagon
% Stable identity: aops-topic-1548800-q01
The shortest closed curve enclosing the circles is the boundary of their convex hull, a ``stretched rubber band''. It consists of
\begin{itemize}
\item six straight segments, each tangent to two neighbouring circles and parallel to a side of the hexagon, so each has length $8$;
\item six circular arcs, one around each circle, between the two tangent segments at that circle.
\end{itemize}

\begin{asy}
usepackage("amsmath");
size(6cm);
defaultpen(fontsize(10pt));
real r=8;
pair[] V; for(int i=0;i<6;++i) V.push(8*dir(60*i));
for(int i=0;i<6;++i) draw(circle(V[i],r),gray(0.6));
draw(V[0]--V[1]--V[2]--V[3]--V[4]--V[5]--cycle,dashed);
for(int i=0;i<6;++i){
  pair n=dir(60*i+30);
  draw(V[i]+r*n--V[(i+1)%6]+r*n,red+linewidth(1.1));
  draw(arc(V[i],r,60*i-30,60*i+30),blue+linewidth(1.1));
  dot(V[i]);}
draw(V[0]--V[0]+r*dir(-30),deepgreen);draw(V[0]--V[0]+r*dir(30),deepgreen);
label("$8$",V[0]+0.72*r*dir(-30),dir(-120),deepgreen);
label("$60^\circ$",V[0]+0.32*r*dir(0),E,fontsize(8pt));
draw(arc(V[0],0.25*r,-30,30),deepgreen);
label("$8$",(V[4]+V[5])/2,N);
label("\mbox{}",(min(currentpicture,user=true)+max(currentpicture,user=true))/2,black); // keeps the text after the figure black when LaTeX typesets the labels inline
\end{asy}

At each vertex the two radii to the tangent points are perpendicular to the two neighbouring sides, so the arc's angle is $360^\circ-90^\circ-90^\circ-120^\circ=60^\circ$. The six arcs together make $360^\circ$, a full circle of radius $8$. The length is
\[
6\cdot8+2\pi\cdot8=48+16\pi\approx98.27 ,
\]
which is $\boxed{98}$ to the nearest whole number.
